\documentclass[12pt,a4paper]{report}
\usepackage[spanish]{babel}
\usepackage{xcolor}
\usepackage{graphicx}
\usepackage{fancyhdr}
\usepackage{geometry}
\usepackage{booktabs}
\usepackage{amssymb}
\usepackage{multirow}
\usepackage{hyperref}
\usepackage{array}
\usepackage{colortbl}

% Configuración de márgenes
\geometry{left=2cm, right=2cm, top=2.5cm, bottom=2.5cm}

% Definir colores personalizados
\definecolor{cafeOscuro}{RGB}{139, 69, 19}
\definecolor{verde}{RGB}{60, 179, 113}
\definecolor{grisClaro}{RGB}{240, 240, 240}
\definecolor{azulOscuro}{RGB}{25, 25, 112}

% Configurar encabezado y pie
\pagestyle{fancy}
\fancyhead[L]{\textcolor{cafeOscuro}{\textbf{KairosMix}}}
\fancyhead[R]{\textcolor{cafeOscuro}{PDCA Temporal}}
\fancyfoot[C]{\thepage}
\fancyfoot[R]{\today}

% Título personalizado
\title{
    \textcolor{cafeOscuro}{\Large \textbf{📊 PDCA Temporal}} \\
    \textcolor{verde}{\Large \textbf{Sistema de Gestión de Frutos Secos Premium}} \\
    \vspace{0.5cm}
    \textcolor{cafeOscuro}{\large KairosMix}
}
\author{}
\date{}

\begin{document}

% Portada personalizada
\begin{titlepage}
    \centering
    \vspace*{3cm}
    
    \textcolor{cafeOscuro}{\Huge \textbf{🌰 PDCA TEMPORAL}}
    
    \vspace{1cm}
    \textcolor{verde}{\huge \textbf{KairosMix}}
    
    \vspace{1cm}
    \textcolor{cafeOscuro}{\Large Sistema de Gestión de Frutos Secos Premium}
    
    \vspace{3cm}
    \rule{\textwidth}{2pt}
    
    \vspace{2cm}
    \textbf{\Large Ciclo de Mejora Continua}
    
    \vspace{1cm}
    \textbf{Plan - Do - Check - Act}
    
    \vspace{3cm}
    
    \begin{flushleft}
        \textbf{Tecnologías:}
        \begin{itemize}
            \item React 19.1.0
            \item Vite 7.0.4
            \item Bootstrap 5.3.7
            \item ESLint
        \end{itemize}
    \end{flushleft}
    
    \vfill
    \textcolor{grisClaro}{\textit{Documento de calidad - Aseguramiento de Calidad}}
    
    \vspace{1cm}
    \textcolor{cafeOscuro}{\textbf{\today}}
    
\end{titlepage}

% Tabla de contenidos
\tableofcontents
\newpage

\chapter{Introducción y Objetivo General}

\section{Objetivo General del PDCA}

Mejorar la calidad, funcionalidad y experiencia de usuario del sistema \textcolor{cafeOscuro}{\textbf{KairosMix}} mediante un ciclo de mejora continua estructurado bajo la metodología PDCA (Plan-Do-Check-Act).

\section{Metodología PDCA}

El ciclo PDCA es una metodología de mejora continua que consta de cuatro fases:

\begin{itemize}
    \item[\textcolor{cafeOscuro}{\textbf{Plan}}] Planificación de objetivos, alcance y actividades
    \item[\textcolor{verde}{\textbf{Do}}] Ejecución de las actividades planificadas
    \item[\textcolor{azulOscuro}{\textbf{Check}}] Verificación y evaluación de resultados
    \item[\textcolor{cafeOscuro}{\textbf{Act}}] Acciones de mejora y seguimiento
\end{itemize}

\chapter{PLAN - Planificación}

\section{Objetivos del Ciclo}

Los objetivos principales de este ciclo PDCA son:

\begin{itemize}
    \item Validar la funcionalidad de los módulos implementados
    \item Identificar defectos y áreas de mejora en el código
    \item Asegurar calidad en la interfaz de usuario (UI/UX)
    \item Mejorar la estructura del código y mantenibilidad
    \item Optimizar el rendimiento de la aplicación web
    \item Establecer métricas de calidad para futuras iteraciones
\end{itemize}

\section{Alcance del Proyecto}

El proyecto \textbf{KairosMix} es una aplicación web moderna para la gestión integral de un negocio de frutos secos premium. El alcance de este PDCA incluye los siguientes módulos:

\subsection{Módulo de Gestión de Productos (Implementado)}

\begin{table}[h]
    \centering
    \begin{tabular}{|l|p{8cm}|}
        \hline
        \textbf{Estado} & ✅ Implementado \\
        \hline
        \textbf{Funcionalidades} & 
            \begin{itemize}
                \item Agregar, editar y eliminar productos
                \item Control de stock inteligente con alertas automáticas
                \item Categorización por tipo de fruto seco
                \item Gestión de proveedores y fechas de vencimiento
                \item Búsqueda y filtrado avanzado en tiempo real
            \end{itemize} \\
        \hline
    \end{tabular}
\end{table}

\subsection{Módulos en Desarrollo}

\begin{table}[h]
    \centering
    \begin{tabular}{|l|p{8cm}|}
        \hline
        \textbf{Módulo} & \textbf{Estado} \\
        \hline
        Gestión de Clientes & 🚧 En desarrollo \\
        \hline
        Gestión de Pedidos & 🚧 En desarrollo \\
        \hline
        Mezcla Personalizada & 🚧 En desarrollo \\
        \hline
    \end{tabular}
\end{table}

\section{Roles y Responsabilidades}

\begin{table}[h]
    \centering
    \begin{tabular}{|l|p{10cm}|}
        \hline
        \rowcolor{grisClaro}
        \textbf{Rol} & \textbf{Responsabilidades} \\
        \hline
        \textbf{Product Owner} & 
            \begin{itemize}
                \item Definir requisitos del proyecto
                \item Validar funcionalidad implementada
                \item Priorizar defectos y mejoras
            \end{itemize} \\
        \hline
        \textbf{Dev Team} & 
            \begin{itemize}
                \item Implementar correcciones
                \item Realizar mejoras técnicas
                \item Optimizar código
            \end{itemize} \\
        \hline
        \textbf{QA/Tester} & 
            \begin{itemize}
                \item Ejecutar pruebas funcionales
                \item Reportar defectos encontrados
                \item Validar criterios de aceptación
            \end{itemize} \\
        \hline
        \textbf{Tech Lead} & 
            \begin{itemize}
                \item Supervisar calidad técnica
                \item Revisar arquitectura del código
                \item Definir estándares de desarrollo
            \end{itemize} \\
        \hline
    \end{tabular}
\end{table}

\section{Métricas a Evaluar}

\begin{itemize}
    \item[\textcolor{cafeOscuro}{\textbf{Funcionalidad}}] Porcentaje de funciones operativas sin defectos críticos
    \item[\textcolor{cafeOscuro}{\textbf{Defectos}}] Clasificación en Críticos, Mayores y Menores
    \item[\textcolor{cafeOscuro}{\textbf{Cobertura de Pruebas}}] Porcentaje de componentes testeados
    \item[\textcolor{cafeOscuro}{\textbf{Performance}}] Tiempo de carga y respuesta de acciones
    \item[\textcolor{cafeOscuro}{\textbf{Accesibilidad}}] Validación de estándares WCAG
    \item[\textcolor{cafeOscuro}{\textbf{Calidad de Código}}] Resultado de ESLint y análisis estático
\end{itemize}

\section{Timeline Planificado}

\begin{table}[h]
    \centering
    \begin{tabular}{|l|p{6cm}|c|}
        \hline
        \rowcolor{grisClaro}
        \textbf{Fase} & \textbf{Descripción} & \textbf{Duración} \\
        \hline
        PLAN & Planificación y definición de objetivos & 1 día \\
        \hline
        DO & Ejecución de actividades y desarrollo & 3-4 días \\
        \hline
        CHECK & Verificación y pruebas & 2-3 días \\
        \hline
        ACT & Mejora continua y optimización & 2-3 días \\
        \hline
        \rowcolor{grisClaro}
        \textbf{Total} & \textbf{Ciclo Completo} & \textbf{1-2 semanas} \\
        \hline
    \end{tabular}
\end{table}

\chapter{DO - Ejecución}

\section{Actividades Planificadas}

\subsection{A. Desarrollo e Implementación}

\subsubsection{Completar módulo de Gestión de Clientes}
\begin{itemize}
    \item Formulario de registro de cliente con validación
    \item Base de datos o almacenamiento local (localStorage/IndexedDB)
    \item Sistema de búsqueda y filtrado de clientes
    \item Vista de detalle de cliente
\end{itemize}

\subsubsection{Completar módulo de Gestión de Pedidos}
\begin{itemize}
    \item Interfaz de creación de pedidos
    \item Asociación de clientes a pedidos
    \item Cálculo automático de totales y costos
    \item Historial y consulta de pedidos anteriores
\end{itemize}

\subsubsection{Implementar Mezcla Personalizada}
\begin{itemize}
    \item Selector interactivo de productos
    \item Cálculo y visualización de información nutricional
    \item Guardado y recuperación de mezclas personalizadas
\end{itemize}

\subsection{B. Testing Inicial}
\begin{itemize}
    \item Pruebas funcionales básicas en cada módulo
    \item Validación de formularios y datos de entrada
    \item Pruebas de almacenamiento y persistencia de datos
    \item Pruebas de integración entre módulos
\end{itemize}

\subsection{C. Optimizaciones Técnicas}
\begin{itemize}
    \item Ejecutar análisis de código: \texttt{npm run lint}
    \item Revisar estructura de componentes React
    \item Optimizar imports y dependencias
    \item Mejorar performance de renderizado
\end{itemize}

\section{Checklist de Implementación}

\textcolor{verde}{\textbf{Antes de pasar a la fase CHECK, validar:}}

\begin{table}[h]
    \centering
    \begin{tabular}{|p{8cm}|c|}
        \hline
        \rowcolor{grisClaro}
        \textbf{Item} & \textbf{Estado} \\
        \hline
        Todos los módulos funcionan en el navegador & ☐ \\
        \hline
        No hay errores en la consola del navegador & ☐ \\
        \hline
        ESLint no reporta errores críticos & ☐ \\
        \hline
        Los datos persisten correctamente & ☐ \\
        \hline
        Responsive design funciona en móvil y desktop & ☐ \\
        \hline
        Aplicación puede ser compilada sin errores & ☐ \\
        \hline
        Documentación básica completada & ☐ \\
        \hline
    \end{tabular}
\end{table}

\chapter{CHECK - Verificación}

\section{Plan de Testing}

\subsection{Pruebas Funcionales}

\subsubsection{Módulo de Productos}
\begin{itemize}
    \item Crear un producto nuevo con validación de campos
    \item Editar información de un producto existente
    \item Eliminar un producto del sistema
    \item Aplicar filtros de búsqueda en tiempo real
    \item Validar alertas automáticas de stock bajo
\end{itemize}

\subsubsection{Módulo de Clientes}
\begin{itemize}
    \item Registrar nuevo cliente con datos completos
    \item Editar información de cliente
    \item Buscar cliente por nombre o email
    \item Validar información obligatoria
    \item Eliminar cliente del sistema
\end{itemize}

\subsubsection{Módulo de Pedidos}
\begin{itemize}
    \item Crear nuevo pedido asociado a un cliente
    \item Agregar productos a un pedido
    \item Calcular total correctamente (incluyen impuestos si aplica)
    \item Generar reporte de pedidos
    \item Editar y cancelar pedidos
\end{itemize}

\subsubsection{Mezcla Personalizada}
\begin{itemize}
    \item Seleccionar múltiples productos
    \item Ver información nutricional combinada
    \item Guardar mezcla personalizada
    \item Recuperar mezclas guardadas previamente
    \item Editar mezclas existentes
\end{itemize}

\subsection{Pruebas de Calidad}

\subsubsection{UI/UX}
\begin{table}[h]
    \centering
    \begin{tabular}{|l|c|}
        \hline
        \textbf{Criterio} & \textbf{Validado} \\
        \hline
        Consistencia visual en toda la aplicación & ☐ \\
        \hline
        Paleta de colores (café y verde) coherente & ☐ \\
        \hline
        Accesibilidad: contraste de colores (WCAG) & ☐ \\
        \hline
        Botones e iconos bien identificados & ☐ \\
        \hline
        Mensajes de error claros y útiles & ☐ \\
        \hline
        Responsive en dispositivos móviles & ☐ \\
        \hline
        Formularios intuitivos y fáciles de usar & ☐ \\
        \hline
    \end{tabular}
\end{table}

\subsubsection{Rendimiento}
\begin{itemize}
    \item Tiempo de carga inicial $< 3$ segundos
    \item No hay lag perceptible en interacciones
    \item Gestión eficiente de memoria (sin memory leaks)
    \item Aplicación funciona bien en conexiones lentas
\end{itemize}

\subsubsection{Código}
\begin{table}[h]
    \centering
    \begin{tabular}{|l|c|}
        \hline
        \textbf{Criterio} & \textbf{Cumple} \\
        \hline
        Sin errores de ESLint & ☐ \\
        \hline
        Componentes reutilizables & ☐ \\
        \hline
        Props validadas correctamente & ☐ \\
        \hline
        Estructura clara de carpetas & ☐ \\
        \hline
        Nombres de variables descriptivos & ☐ \\
        \hline
        Código comentado donde sea necesario & ☐ \\
        \hline
    \end{tabular}
\end{table}

\section{Reporte de Defectos}

\subsection{Plantilla de Defecto}

\begin{table}[h]
    \centering
    \begin{tabular}{|p{3cm}|p{10cm}|}
        \hline
        \textbf{Campo} & \textbf{Descripción} \\
        \hline
        \textbf{ID} & Código único del defecto (ej: DEFECTO-001) \\
        \hline
        \textbf{Severidad} & Crítico | Mayor | Menor \\
        \hline
        \textbf{Módulo} & Nombre del módulo afectado \\
        \hline
        \textbf{Descripción} & Descripción clara del problema \\
        \hline
        \textbf{Pasos para reproducir} & Lista detallada 1. 2. 3. \\
        \hline
        \textbf{Resultado esperado} & Comportamiento correcto \\
        \hline
        \textbf{Resultado actual} & Comportamiento observado \\
        \hline
        \textbf{Asignado a} & Miembro responsable \\
        \hline
        \textbf{Estado} & Abierto | En progreso | Resuelto \\
        \hline
    \end{tabular}
\end{table}

\section{Matriz de Evaluación}

\begin{table}[h]
    \centering
    \begin{tabular}{|l|c|c|p{6cm}|}
        \hline
        \rowcolor{grisClaro}
        \textbf{Criterio} & \textbf{Peso} & \textbf{Estado} & \textbf{Observaciones} \\
        \hline
        Funcionalidad & 40\% & ⭕ & \\
        \hline
        Calidad de Código & 25\% & ⭕ & \\
        \hline
        UX/UI & 20\% & ⭕ & \\
        \hline
        Rendimiento & 10\% & ⭕ & \\
        \hline
        Testing Coverage & 5\% & ⭕ & \\
        \hline
        \rowcolor{grisClaro}
        \textbf{TOTAL} & \textbf{100\%} & \textbf{⭕} & \\
        \hline
    \end{tabular}
\end{table}

\textbf{Leyenda:} 🟢 Aprobado $|$ 🟡 Parcial $|$ 🔴 No aprobado $|$ ⭕ Por evaluar

\section{Resultados Esperados}

Los criterios de aceptación para esta fase CHECK son:

\begin{itemize}
    \item Mínimo 85\% de funcionalidad operativa
    \item 0 defectos críticos identificados
    \item Máximo 5 defectos mayores
    \item ESLint sin errores o solo warnings menores
    \item Score de accesibilidad $\geq 80$
    \item Tiempo de carga $< 3$ segundos
\end{itemize}

\chapter{ACT - Acciones de Mejora}

\section{Análisis de Resultados}

\subsection{Actividades}
\begin{itemize}
    \item Documentar hallazgos principales del ciclo
    \item Priorizar defectos encontrados por severidad
    \item Identificar raíces de problemas comunes
    \item Clasificar mejoras por impacto y complejidad
    \item Generar lecciones aprendidas
\end{itemize}

\section{Plan de Mejora Continua}

\subsection{Mejoras Inmediatas (Sprint Siguiente)}

\begin{itemize}
    \item Corregir todos los defectos críticos identificados
    \item Optimizar componentes con bajo rendimiento
    \item Mejorar mensajes de error y validaciones de formularios
    \item Implementar confirmaciones para acciones destructivas
\end{itemize}

\subsection{Mejoras a Corto Plazo (2-4 semanas)}

\begin{itemize}
    \item Implementar pruebas unitarias con Jest
    \item Agregar logging y monitoreo en consola
    \item Documentar API interna del proyecto
    \item Mejorar accesibilidad a nivel WCAG AA
    \item Crear guía de estilo de componentes
\end{itemize}

\subsection{Mejoras a Largo Plazo (1-2 meses)}

\begin{itemize}
    \item Integración con backend real (Node.js/Express)
    \item Sistema de autenticación robusto con JWT
    \item Base de datos SQL o NoSQL
    \item Análisis y reportes avanzados
    \item Exportación a Excel mejorada con formatos
    \item Sistema de notificaciones en tiempo real
\end{itemize}

\section{Acciones de Seguimiento}

\begin{table}[h]
    \centering
    \begin{tabular}{|p{3cm}|p{3cm}|p{2cm}|p{3cm}|c|}
        \hline
        \rowcolor{grisClaro}
        \textbf{Acción} & \textbf{Responsable} & \textbf{Fecha} & \textbf{Dependencias} & \textbf{Estado} \\
        \hline
        & & & & ☐ \\
        \hline
        & & & & ☐ \\
        \hline
        & & & & ☐ \\
        \hline
    \end{tabular}
\end{table}

\section{Plan para Próximo Ciclo PDCA}

\begin{itemize}
    \item Definir nuevos objetivos de mejora basados en resultados
    \item Ajustar métricas de evaluación si es necesario
    \item Incorporar feedback de usuarios finales
    \item Planificar incremento de funcionalidades
    \item Mantener registro histórico de mejoras
\end{itemize}

\chapter{Indicadores Clave (KPIs)}

\section{Seguimiento de Métricas}

\begin{table}[h]
    \centering
    \begin{tabular}{|l|c|c|c|}
        \hline
        \rowcolor{grisClaro}
        \textbf{KPI} & \textbf{Meta} & \textbf{Actual} & \textbf{Estado} \\
        \hline
        \% Funcionalidad Operativa & 85\% & ⭕ & \\
        \hline
        Defectos Críticos & 0 & ⭕ & \\
        \hline
        Defectos Mayores & $\leq 5$ & ⭕ & \\
        \hline
        Tiempo Promedio de Carga & $< 3s$ & ⭕ & \\
        \hline
        Cobertura de Testing & $\geq 50\%$ & ⭕ & \\
        \hline
        Score ESLint & 100\% & ⭕ & \\
        \hline
        Accesibilidad (WCAG) & $\geq 80$ & ⭕ & \\
        \hline
    \end{tabular}
\end{table}

\chapter{Información del Equipo}

\section{Miembros del Proyecto}

\begin{table}[h]
    \centering
    \begin{tabular}{|l|l|l|}
        \hline
        \rowcolor{grisClaro}
        \textbf{Nombre} & \textbf{Rol} & \textbf{Email} \\
        \hline
        & Product Owner & \\
        \hline
        & Tech Lead & \\
        \hline
        & Developer 1 & \\
        \hline
        & Developer 2 & \\
        \hline
        & QA/Tester & \\
        \hline
    \end{tabular}
\end{table}

\section{Escalaciones y Contacto}

\begin{itemize}
    \item[\textbf{Problemas Técnicos}] Contactar al Tech Lead
    \item[\textbf{Defectos Críticos}] Revisar e implementar corrección inmediata
    \item[\textbf{Cambios de Scope}] Consultar con Product Owner y replanificar
    \item[\textbf{Bloqueos}] Escalar a Tech Lead para desbloqueo urgente
\end{itemize}

\chapter{Recursos Técnicos}

\section{Stack Tecnológico}

\begin{table}[h]
    \centering
    \begin{tabular}{|l|l|}
        \hline
        \rowcolor{grisClaro}
        \textbf{Componente} & \textbf{Tecnología} \\
        \hline
        Frontend Framework & React 19.1.0 \\
        \hline
        Build Tool & Vite 7.0.4 \\
        \hline
        CSS Framework & Bootstrap 5.3.7 \\
        \hline
        Iconografía & Lucide React \\
        \hline
        Routing & React Router v7 \\
        \hline
        Exportación & XLSX (Excel) \\
        \hline
        Alerts & SweetAlert2 11.22.2 \\
        \hline
        Linting & ESLint 9.30.1 \\
        \hline
    \end{tabular}
\end{table}

\section{Documentación de Referencia}

\begin{itemize}
    \item \href{https://react.dev}{Documentación Oficial React}
    \item \href{https://vitejs.dev}{Documentación Vite}
    \item \href{https://getbootstrap.com/docs/5.0}{Bootstrap 5 Documentation}
    \item \href{https://eslint.org/docs/rules}{ESLint Rules}
    \item \href{https://www.w3.org/WAI/WCAG21/quickref}{WCAG 2.1 Accessibility Guidelines}
\end{itemize}

\chapter*{Registro de Cambios}

\section*{Ciclo \#1}

\begin{table}[h]
    \centering
    \begin{tabular}{|l|c|p{8cm}|}
        \hline
        \rowcolor{grisClaro}
        \textbf{Fase} & \textbf{Resultado} & \textbf{Notas} \\
        \hline
        PLAN & ✅ Completado & Objetivos y alcance claros definidos \\
        \hline
        DO & 🔄 En Progreso & Implementación en curso \\
        \hline
        CHECK & ⭕ Pendiente & Próxima fase a ejecutar \\
        \hline
        ACT & ⭕ Pendiente & Acciones de mejora pendientes \\
        \hline
    \end{tabular}
\end{table}

\chapter*{Conclusión}

Este documento \textcolor{cafeOscuro}{\textbf{PDCA Temporal}} proporciona un marco estructurado para la mejora continua del proyecto \textbf{KairosMix}. Mediante la aplicación rigurosa de las fases Plan-Do-Check-Act, el equipo puede:

\begin{itemize}
    \item ✓ Identificar y resolver defectos de manera sistemática
    \item ✓ Establecer estándares de calidad claros
    \item ✓ Mejorar continuamente la arquitectura y código
    \item ✓ Optimizar la experiencia del usuario
    \item ✓ Documentar lecciones aprendidas
    \item ✓ Preparar la base para ciclos PDCA futuros
\end{itemize}

\vspace{1cm}

\textcolor{cafeOscuro}{\textbf{Última actualización:}} \today

\textcolor{cafeOscuro}{\textbf{Próxima revisión:}} Una semana después de iniciado el ciclo

\vspace{1cm}

\textit{\textcolor{grisClaro}{Este es un documento dinámico de Aseguramiento de Calidad que debe ser actualizado continuamente según los avances del proyecto.}}

\end{document}
